\section{Persistence in commerce: the five-year proof of continued use}\label{sec:gate}

This section is about one outcome: whether a registered mark's owner
proved, between the fifth and sixth year of the registration, that the
product the mark names was still in commerce. The sample is every unique
registration issued 2002--2018 that carries a score; failure is a dated
cancellation for non-use at registration age 4.0--8.5 years
(Section~\ref{sec:gates}). Half of registered marks fail it. How the
penalty found here moved across industries and decades is
Section~\ref{sec:waves}; what became of the firms is
Section~\ref{sec:success}.

\subsection{Why the analysis conditions on registration}\label{sec:filter}

The question is about companies whose products entered commerce: given
that a product was brought to market, does the language of its first
description predict whether it was still there five years on? Registration
is the record's marker of entry --- a use-based application asserts the
mark is in commerce, and an intent-to-use application registers only after
a sworn statement of use --- so the analysis begins there. An application
that never registered left the record, nine times in ten, by the
applicant's own inaction, before any product met a customer. Including
those applications would change the question into one about everyone who
considered launching, including those who then did not, and since they are
2.3 million of the 6.0 million scored applications, their inclusion also
buries the commercial signal in launch noise.
Section~\ref{sec:robust_uncond} shows the pooled analysis for a reader
skeptical of the conditioning: the registration step's shape is stamped
onto every later outcome, and the gradients found below are attenuated and
in places reversed.

Conditioning would miscount if abandonment were usually a pause rather
than an ending --- if owners reworked refused applications and refiled,
the same product would be counted once as a failure and again as a
success. It is not: among 2.28 million abandoned applications, 3.6\% are
refiled within six years and 1.5\% lead to a registration
(Section~\ref{sec:amend}). Crediting every successful refiling as an
eventual registration moves the registration profile by less than a tenth
of a point. The refiled cases are $1.1$\% of this section's sample, they
fail the five-year proof at nearly the sample rate (51.7\% against
52.8\%), and the headline contrast below is $+2.25$pp with them and
without them.

\subsection{Leading marks fail the five-year proof more often}\label{sec:gate_headline}

Marks written in language their industry subsequently moved toward fail
the five-year proof more often than marks written in language it was
moving away from. Within Nice class and registration year, controlling for
whether the filer drafted the mark themselves or hired counsel, and for
the length of the description, registrations in the most leading fifth of
lead fail $1.4$ points more often (53.3\%) than those in the most lagging
fifth (51.9\%; $t = 8.3$, Table~\ref{tab:gate_lpm}).

Among 4.12 million registration class-records 2002--2018, dated
cancellation at the five-year proof rises from 50.0\% at the bottom
lead quintile to 52.4\% at the top, positive in 33 of 44 classes with
sufficient volume (Figure~\ref{fig:forest}). But a raw pooled contrast
mixes the within-industry gradient with class and registration-year
composition, treats multi-class registrations as independent observations,
and ignores that outcomes correlate within owner. The primary estimates
therefore come from a linear probability model on 3.10 million unique
registrations, one row per serial, with lead quintiles cut \emph{within}
Nice class and registration year. The specification adds
class$\times$registration-year fixed effects; controls for counsel of
record, intent-to-use basis, log description length, log owner filing
volume, owner domicile, and foreign filing basis;\footnote{A US
application is filed either on \emph{use} --- the mark is already in
commerce --- or on \emph{intent to use}, which defers proof until after
allowance. Foreign applicants may instead file on a home registration
(\S44(e)) or extend an international registration into the US under the
Madrid Protocol (\S66(a)); the latter maintain under \S71 rather than
\S8. These are statutory routes to the same registration, and they differ
sharply in who files and how often the mark is abandoned, which is why
each enters as a control.} and standard errors clustered on 1.27 million
normalized owners (Table~\ref{tab:gate_lpm}).\footnote{Tied scores are common (17.1\% of registrations share a lead
value within their cell); tie-breaking is deterministic and immaterial,
moving the contrast by at most $0.002$pp against $2.287$ (supplementary
material).}

\begin{table}[h]\centering\small
\caption{Failure at the five-year proof of continued use and lead:
design-based estimates. LPM of event-dated failure on
within-class$\times$registration-year lead quintiles (Q1 omitted), unique
serials, registrations 2002--2018. Failure is a dated cancellation for
non-use at registration age 4.0--8.5 years, under \S8 or, for Madrid
\S66(a) registrations, \S71. All specifications include
class$\times$registration-year fixed effects; controls are counsel of
record, intent-to-use basis, log goods/services length, log owner filing
count, owner domicile (US, CN), and foreign basis (\S44(e), \S66(a)).
Base failure rate 52.3\% on this sample.}
\label{tab:gate_lpm}
\begin{tabular}{llrrr}
\toprule
Sample & Specification & Q5$-$Q1 (pp) & SE & $n$ \\
\midrule
All & FE only & $+2.29$ & $0.16$ & 3{,}104{,}534 \\
All & FE + controls & $+1.39$ & $0.17$ & 3{,}104{,}534 \\
Counsel-represented & FE + controls & $+1.15$ & $0.20$ & 2{,}414{,}318 \\
Counsel + US-domiciled & FE + controls & $+1.60$ & $0.22$ & 1{,}888{,}352 \\
Self-filed & FE + controls & $+1.80$ & $0.26$ & 690{,}216 \\
Excluding Madrid \S66(a) & FE + controls & $+1.82$ & $0.18$ & 2{,}924{,}427 \\
\midrule
Counsel + US, 2002--07 & FE + controls & $-0.25$ & $0.33$ & 551{,}989 \\
Counsel + US, 2008--12 & FE + controls & $+2.00$ & $0.42$ & 551{,}924 \\
Counsel + US, 2013--18 & FE + controls & $+2.65$ & $0.25$ & 784{,}439 \\
\midrule
All & continuous, pp per sd of lead & $+0.58$ & $0.05$ & 3{,}104{,}534 \\
Counsel + US & continuous, pp per sd of lead & $+0.63$ & $0.06$ & 1{,}888{,}352 \\
\bottomrule
\end{tabular}
\\[0.3em]\footnotesize Standard errors replace significance stars: with
$n$ in the millions every coefficient has $p < 0.01$, so stars would mark
each row identically. Errors are clustered on \emph{normalized owner} ---
names uppercased, stripped of punctuation and of two dozen legal suffixes,
so \textsc{acme corp.} and \textsc{acme corporation} are one cluster. A
single owner files many marks, drafts them in a house style, and abandons
them together when the business fails; the within-owner correlation of
outcomes at the five-year proof is $0.38$ pooled across classes ($0.42$ in
the single-class checks where it was first measured), so the 1.27 million
owner clusters, not the filings, carry the information. The Madrid row
drops \S66(a) registrations rather than absorbing them in a control, since
they maintain under a different statute; the estimate is larger without
them.
\end{table}

The gradient is monotone within cells. Whether the filing was drafted by
counsel or by the owner is held because it moves both sides of the
comparison at once: self-filers fail the five-year proof at 69.6\% against
46.8\% for represented owners, and Section~\ref{sec:representation} showed
they write shorter, more class-typical, slightly more leading
descriptions, so without the control part of any lead gradient would be a
self-filing gradient. With it, the penalty is of the same order in both
strata ($+1.8$pp self-filed, $+1.2$pp counselled). The full control set
cuts the magnitude by about a third without absorbing it, and the estimate
is if anything \emph{largest} in the cleanest stratum ---
counsel-represented, US-domiciled owners --- so it is not carried by the
self-filed or foreign mass-filing segments whose failure rates are high
for their own reasons.\footnote{The level effects on those covariates are
large in their own right: counsel $-19.5$pp and Madrid \S66(a) basis
$+10.3$pp.} About forty percent of the raw pooled contrast is composition:
$+1.4$pp design-based against $+2.5$pp raw. The penalty is also not
constant over the sample: within counsel-represented US-domiciled owners
the earliest third of cohorts carries none of it ($-0.25$pp, SE $0.33$)
--- leading marks registered 2002--2007 were no less likely to survive
than lagging ones --- while the later thirds carry $+2.00$ and $+2.65$pp.
Section~\ref{sec:waves} takes that variation as its subject.

\begin{figure}[h]\centering
\includegraphics[width=0.72\textwidth]{results/event_gate_forest.png}
\caption{Failure lift at the five-year proof of continued use (lead top
minus bottom quintile) by industry, registrations 2002--2018, raw quintile
contrasts; classes need 5{,}000 scored registrations to enter. Positive
bars indicate the leading penalty. Whiskers are two-sample binomial
intervals, uncorrected for the within-owner correlation ($0.38$) and so
narrower than clustered intervals; read Table~\ref{tab:gate_lpm} for
precision. Per-class breakouts are in the online appendix.}
\label{fig:forest}
\end{figure}

The per-class extremes of the raw screen are in the supplementary
material and the online appendix.

The penalty is a rate, not a timing, and not disengagement. The proof is
close to a fixed appointment (96.1\% of cancellations post at
registration age 6.5--7.0), so a mark failing outside the 4.0--8.5 window
would count as a survivor; on the 3.15 million registrations whose ninth
anniversary is inside the record, where failure is settled and no window
is needed, the contrast is $+2.71$pp against the window's $+2.65$, and
conditional on failing, failing late carries a contrast of $-0.01$pp
($t = -0.2$). Nor is the proof measuring an owner who lost interest:
passage is flat across office-action response-speed quintiles (46.1\% to
48.1\%, medians seven days to 183) and the penalty is present in both
speed halves ($-3.48$ and $-2.36$pp on passing, against $-2.94$ pooled).
The supplementary material carries the detail, including why 2019 is the
last observable cohort.

\subsection{The penalty under three theme scorings}\label{sec:three_scorings}

Every estimate above scored filings on fifty themes fitted once over the
whole corpus, with each class supplying its own reference. That is one of
three ways to partition the vocabulary, and the penalty does not depend on
the choice. The identical registrations were re-scored under
\emph{500 global themes} (the same construction at ten times the
resolution, so a theme is closer to a product category than to an
industry) and under \emph{50 themes fitted per class} (a separate model on
each class's own filings, so nothing is shared across classes). In all
three, surprise is measured against the filing's own class on per-filing
windows of 1{,}826 days each side.

\begin{figure}[h]\centering
\includegraphics[width=\textwidth]{results/fig_resolution_compare.png}
\caption{Failure at the five-year proof of continued use by quintile of
lead (left) and of atypicality (right) under the three theme scorings, on
the identical 3.10 million registrations 2002--2018; quintiles cut within
Nice class and registration year, no controls. Whiskers are 95\% binomial
intervals per quintile, uncorrected for owner clustering. The lead
gradient survives every scoring; the atypicality gradient does not
survive per-class themes.}
\label{fig:resolution}
\end{figure}

The lead penalty is a property of the filing, not of the partition
(Figure~\ref{fig:resolution}, left): $+2.29$pp under the production
scoring, $+2.61$pp at 500 themes, and $+1.94$pp under per-class themes,
each with a standard error of $0.09$, and positive under all three at once
in 28 of the 44 classes with at least 3{,}000 sampled registrations. The
scorings agree only moderately filing by filing --- lead correlates at
$0.64$ between the two global resolutions and at $0.52$ between the global
and the per-class fifty --- so a third to a half of the variance in any
one scoring's lead is particular to that scoring, and the penalty survives
the disagreement. (Table~\ref{tab:resolution} tabulates the contrasts; the
standard errors here are two-sample binomial and unclustered, and
outcomes correlate within owner at $0.38$, so single-class $t$-statistics
overstate precision by roughly a factor of two.)

\begin{table}[h]\centering\small
\caption{Failure contrasts at the five-year proof under three theme
scorings, on the identical registrations. Top-minus-bottom quintile,
quintiles cut within Nice class and registration year, no controls;
$n = 3{,}105{,}048$, base failure 52.3\%.}
\label{tab:resolution}
\input{results/resolution_compare.tex}
\end{table}

Atypicality behaves differently (Figure~\ref{fig:resolution}, right).
Under the two global scorings the most atypical fifth of registrations
fails the five-year proof $4.3$--$4.4$ points less often than the least
atypical; under per-class themes the gap is $0.7$ points. The two global
scorings' atypicality correlates at $0.79$; global and per-class
atypicality correlate at $0.39$. A global model places every class in one
set of coordinates, so a filing that imports a theme the class rarely uses
and other classes use often registers as atypical for its class. A model
fitted on the class alone re-partitions the class's own vocabulary, and
the imported theme becomes a local theme like any other. What the global
scoring calls atypical is, in large part, drawing on a theme from
elsewhere.

That population can be measured directly, and it separates two kinds of
new. Under the 500-theme model, for each theme and class, the first
sustained year is the first year by which the theme has been the dominant
theme of at least one hundred filings in that class; the corpus-wide first
sustained year is the earliest of those across classes.\footnote{Sustained
years and theme shares are computed on a 25\% systematic sample of all
filings (every fourth serial), scaled up, over filing years 1986--2024.}
Each registration in the sample is assigned its dominant theme and
classified at its filing date as carrying a theme new to the corpus, a
theme new to its class (established elsewhere, arriving here), or an
established theme. Of 4.12 million registration class-records, 3.97
million carry an established theme, 151{,}012 a theme new to their class,
and 779 a theme new to the corpus --- the last too few to support a
claim, because almost every theme a 500-theme model can resolve had been
established in some class before 2002.

Registrations carrying a theme new to their class fail the five-year proof
at 48.3\% against 51.6\% for established themes, and $0.43$ points less
often within class and registration year (SE $0.13$). Holding the
filing's own lead and atypicality as well reverses the sign: net of them,
the new-to-class registrations fail $1.06$ points more often. The two
statements describe one population. A theme arrives in a class in filings
that are atypical for that class, and atypical filings survive; that is
the $-0.43$. Among filings of equal atypicality, the ones carrying the
arriving theme fail more; that is the $+1.06$. The atypicality protection
of the global scorings is therefore mostly the protection of extension:
language established elsewhere and extended into a new class survives.
The penalty attaches to being early with language, not to being far from
the class's centre.

Nor is the lead penalty an atypicality effect in disguise. The two could
be confounded mechanically --- lead is the difference of two nonnegative
levels whose average is atypicality, so its spread grows with atypicality
and its magnitude correlates with it at $+0.29$ --- but the signed value
does not ($-0.05$), and cutting lead quintiles within fifths of
atypicality leaves the penalty intact in every fifth: $+2.1$, $+3.1$,
$+3.3$, $+1.7$ and $+2.5$pp from the least atypical fifth to the most
(each SE under $0.3$).

\subsection{The risk persists at the year-ten renewal}\label{sec:renewal}

Survivors of the five-year proof face a combined declaration and renewal
at year ten, and on the design-based estimates the lead penalty does not
attenuate there: $-0.49$pp per standard deviation on passing the
five-year proof, $-0.63$pp on renewing at ten (Table~\ref{tab:staged},
both stated on survival). Pooled across classes the renewal profile is a
V rather than a slope --- renewal is highest just below the centre of the
lead axis and lowest at both tails, the lagging tail lowest of all
(45.5\% against 50.6\% at the most leading; fit detail in the
supplementary material). Most of the difference between leading and
lagging resolves at the five-year proof; what remains at renewal is a
penalty on both extremes, the lagging one the larger.

\subsection{What the five-year proof shows}\label{sec:gate_conclusion}

Among products that entered commerce, being described in language the
industry was moving toward predicts being gone within six years: $+1.4$
points on a 52.3\% base with the full design, $+2.3$ raw, in three of
every four classes, under every partition of the vocabulary, robust to
timing, to drafting, and to who files. Atypical language predicts the
opposite, but that protection belongs mostly to filings extending a theme
established elsewhere into a new class; net of atypicality, even those
arriving-theme filings fail more. The market's step and the examiner's
step read the text differently for their own reasons: the examiner reads
distance from the class's usual language at a moment when lead does not
yet exist to be read (Section~\ref{sec:reg_rates}); the five-year proof is
recorded after the industry's language has moved, and it is where being
early shows up as a cost.
